\documentclass[12pt]{article}
\usepackage[utf8]{inputenc}
\usepackage[T1]{fontenc}
\usepackage{hyperref}
\title{M\'enik Language Resources for K-12 Curriculum Development}
\author{Compiled by Claude Sonnet 5 (claude-sonnet-5) for Andr\'as Kornai}
\date{}
\begin{document}
\maketitle

\section{Introduction}
M\'enik (ISO 639-3 tnr) is also known as Bedik, Budik or Tandanke. It is spoken in the Bandafassi area of eastern Senegal, near the Bassari. Sources give about 3,500 people. The ELDP describes many speakers as using it as a second language or having stopped using it. Some speakers mix it heavily with Pulaar and Bassari. One search result lists it under Mande, which conflicts with the Atlantic-Congo classification in the prompt.

\section{Primary Educational Resources}
No primers or school materials were found in the search results.

\section{Linguistic and Reference Materials}
The ELDP project ``Documentation of Bedik'' (collection 0141, Adjaratou Oumar Sall, IFAN, Dakar) is described as the first description of the language. It includes annotated texts and a first dictionary of about 3,000 entries. I did not verify its contents or access conditions. A formal grammar was not found.
\begin{itemize}
\item \url{https://eldp.access.preservica.com/dk0141/}
\end{itemize}

\section{Religious and Community Resources}
No scripture or community organisation was found.

\section{Digital and Audio Resources}
\begin{itemize}
\item ELDP/ELAR recordings and texts, as above: \url{https://eldp.access.preservica.com/dk0141/}
\item Endangered Languages Project entry: \url{https://www.endangeredlanguages.com/node/110130}
\end{itemize}

\section{Curriculum Development Context}
Only documentation material was found, and no teaching materials. A curriculum would start from the ELDP collection and work with the community, whose language use is mixed. This dossier is short because little was found.
\end{document}